\paragraph{Jump-enabled multi-asset comparison.}

To test whether the temporal benefit of the single-asset jump
mechanism survived composition, we compared naive and hybrid
composition under three jump settings (Table~\ref{tab:jump_ablation}).
Unlike the six-method comparison, which used observed SPY in
training, this comparison simulated the SPY market path in both the
$2{,}766$-day training window and the $249$-day holdout, so a jump
setting could apply to the market generator. Jumps were disabled,
enabled only in the SPY market generator, or enabled in both SPY and
every asset generator. We froze the closing-price hidden Markov generators and SIM
calibrations of the six-method comparison. Enabled models
received the SPY settings of Table~\ref{tab:hyperparams},
$(\epsilon,\lambda)=(10^{-4},90)$, $p_{\rm neg}=0.52$, and
$N_{\rm tail}=5$, without per-asset tuning. These settings were tuned on volume-weighted SPY prices and applied
here to generators fitted to closing prices. At these settings,
most paths contain no episode, especially over the $249$-day holdout,
so the jump effects are carried by a minority of paths. In both windows and for all three jump settings, the correction used
each simulated market path's variance and each asset draw's variance.
This differs from the six-method holdout, which fixed the market
variance at its training value.

For each jump setting, naive and hybrid composition received the same
asset and market draws, and one market path was
shared across all assets in each replication. We generated $250$
paths for each of four seeds ($1234$, $2345$, $3456$, and $4567$),
giving $1{,}000$ replications per asset, setting, and method. We used a path-level version of ACF-MAE, computed on each path and
then averaged across paths (Eq.~\eqref{eq:acf-path-error} in Online
Appendix~\ref{sec:supp-validation}); the single-asset ACF-MAE instead
averages the path ACFs before measuring the error. We used
lags $1$--$25$, the lag window of the tuning objective, with lags
$1$--$60$ as a secondary check, and repeated the marginal and variance
diagnostics. The temporal scores of the single-asset and jump comparisons
therefore have different estimands and cannot be compared directly. We included all paths, not only jump-active ones, in the summaries;
jump-active subsets only described the mechanism
(Table~\ref{tab:jump_ablation_episodes}). For Monte Carlo uncertainty, we averaged each metric across tickers
within each replication, which kept assets together under their
shared market path. Jump settings were paired by replication, sharing
the draws of any component whose jump setting did not change, and the
approximate $95\%$ confidence intervals were the mean of the $1{,}000$
replication-level differences plus or minus $1.96$ standard errors
(Table~\ref{tab:jump_ablation_uncertainty}). These intervals describe simulation
uncertainty conditional on the fitted models and observed histories,
not uncertainty across market regimes.
